\chapter{Data Conversion Service Specification}
\section*{\textit{Version 1.0}}
\section{Introduction}

A conscious design decision was made for CIShell not to enforce a central data
model/format that all algorithms had to work with. Instead, an algorithm
expresses in its service metadata for each data item coming into and out of the
algorithm what the format of the data is. It is the job of the code calling the
algorithm to get the data in the right format before calling the algorithm. The
Data Conversion Service is used here to simplify the process of converting data.

\subsection{Entities}

\begin{itemize}
  \item \textit{DataConversionService} - The service interface for converting
  data to different formats.
  \item \textit{Converter Algorithm} - A special type of algorithm which
  converts data from one format to another.
  \item \textit{Converter} - The interface for a wrapped set of converter
  algorithms returned by the \class{DataConversionService} that will convert
  data from one format to another.
  \item \textit{Validator} - A special type of converter algorithm which checks
  a specified file to see if it is a particular file format. 
\end{itemize}

\section{Data Conversion Service}

The Data Conversion Service provides unified access to converter algorithms.
\class{DataConversionService} system developers may choose not to use converter
algorithms, but is ill-advised. Also, good implementations will take advantage of
the nature of algorithms to allow for more than just single hop conversions.
Since all converter algorithms specify a single data object in and a single data
object out, a graph can be constructed where nodes are the data formats and edges
are the converters. Using this directed graph, when a conversion between data
formats is requested, the \class{DataConversionService} will choose the shortest
path of converters to do the conversion. A hypothetical conversion graph is
illustrated in figure \ref{fig:conversionGraph}.

\begin{figure}[htb!]
\centering
\includegraphics[width=70mm]{../img/conversionGraph.pdf}
\caption{A Conversion Graph}
\label{fig:conversionGraph}
\end{figure}

\subsection{Converter Algorithms}

A converter algorithm is a standard CIShell algorithm with the following
constraints: 

\subsubsection*{Required:}
\begin{itemize}
  \item The algorithm must be a conformant \class{AlgorithmFactory}
  implementation and properly registered as a service.
  \item The algorithm must take in a single \class{Data} item and convert the
  item producing a single \class{Data} item. This must be reflected in the
  algorithm's service metadata where ``in\_data'' and ``out\_data'' have only
  one format each.
  \item The algorithm's service metadata must contain a legal ``service.pid''.
  \item The algorithm's service metadata must have ``type=converter''.
  \item The algorithm's service metadata must have ``conversion=lossy'' if
  data is lossed during conversion or ``conversion=lossless'' if not.
\end{itemize}

\subsubsection*{Optional:}
\begin{itemize}
  \item The algorithm's service metadata should have ``remoteable=true'' if it
  meets the requirements of a remoteable algorithm.
  \item The algorithm's service metadata should have a ``label'' which is a
  short human-readable name for the converter, usually with a readable name of
  the input data format and output data format.
  \item The algorithm's service metadata should have a ``description''
  describing the conversion in more detail, perhaps going into what data may be
  lossed if ``conversion=lossy''.
\end{itemize}

\subsection{Validator Algorithms}

A validator algorithm is a standard CIShell algorithm which checks either an
incoming or outgoing file to be sure it is of the type specified. These algorithms
have the following constraints:

\subsubsection*{Required:}
\begin{itemize}
  \item The algorithm must be a conformant \class{AlgorithmFactory}
  implementation and properly registered as a service.
  \item The algorithm must take in a single \class{Data} item and validate the
  item producing a single \class{Data} item (with the same data, but changed
  format) or \class{null} if the file being validated is not of the right
  type. This must be reflected in the algorithm's service metadata where
  ``in\_data'' and ``out\_data'' have only one format each with one containing
  a ``file:'' format and the other a ``file-ext:'' depending on the direction
  of validation.
  \item The algorithm must not alter the data. Its only purpose is to validate
  the proposed incoming or outgoing file.
  \item The algorithm's service metadata must contain a legal ``service.pid''.
  \item The algorithm's service metadata must have ``type=validator''.
\end{itemize}

\subsubsection*{Optional:}
\begin{itemize}
  
  \item The algorithm's service metadata should have ``remoteable=true'' if it
  meets the requirements of a remoteable algorithm.
  \item The algorithm's service metadata should have a ``label'' which is a
  short human-readable name for the converter, usually with a readable name of
  the input data format and output data format.
  \item The algorithm's service metadata should have a ``description''
  describing the conversion in more detail, perhaps going into what data may be
  lossed if ``conversion=lossy''.
\end{itemize}

\orgcishellserviceconversion{}
